\section{本章小结}
\label{sec:mfo-summary}

本章从一个问题出发：可行集是球面或正交矩阵集合时，能否把它本身当作计算空间。答案是肯定的，但要付出代价：约束从问题的表述中消失，转而进入每一步的几何运算。所求解的仍然是约束极值，黎曼梯度为零就是Lagrange乘子驻定条件，切向曲率就是Lagrange函数的Hessian在切空间上的二次型。

前提是可行集具有光滑结构，这由正则等式约束保证。集合、度量和表示是三种不同的选择。集合决定极小点，度量决定梯度和曲率的尺度，表示决定成本。十字形交点、棱角这样的点上没有线性切空间，本章的迭代不能直接使用。

全部方法共用同一个循环：在切空间中产生候选步，回缩到流形，用拉回目标检验，再搬运历史向量。回缩只需要零步不动和一阶一致性，不必是测地线，也不必是最近点；二阶回缩使拉回Hessian等于黎曼Hessian，一般回缩只在驻点处相等。向量搬运只保证落入新切空间并保持线性，内积和长度的保持要逐个核对。

表~\ref{tab:mfo-method-summary}用贯穿例对比各种方法需要的信息及其结果，保证的来源各不相同。梯度下降的残差界依赖拉回的光滑性；共轭梯度和拟Newton方法还依赖线搜索、搬运与回缩的配套；Newton法在非退化解附近快速收敛，但它寻找驻点，远离解时可能上升；信赖域用比值检验曲率模型，需要处理负曲率才能得到二阶保证。

\begin{BookInlineTable}
 \centering
 \caption{本章各方法的比较。}
 \label{tab:mfo-method-summary}
 \begin{tabular}{lll}
  \toprule
  方法 & 循环外还需要 & 贯穿例或实验中的表现\\
  \midrule
  梯度下降 & 线搜索 & 每轮误差约减半\\
  共轭梯度 & 搬运、重启 & 低维无优势，高维更快\\
  Newton法 & Hessian作用 & 近解立方，远离时可能上升\\
  BFGS & 等距搬运、曲率检查 & 逐步学习曲率\\
  信赖域 & 半径、比值检验 & 利用负曲率，四轮收敛\\
  \bottomrule
 \end{tabular}
\end{BookInlineTable}

算法的保证由度量、模型、搬运和误差控制共同决定，欧氏结论不能不加核对地移用。同一个子空间有多种矩阵记录时，商结构去掉冗余的竖直方向，使非退化条件恢复。与一般约束方法比较时，要以同一目标和精度同时报告可行误差与黎曼驻点残差，并把投影、回缩和搬运的代价计算在内。

本章的度量都是诱导度量。概率分布族上的Fisher度量给出另一种几何，对应的自然梯度方向的统计含义可参考\BookChapterRef{chap:information-geometry}{卷七第28章“信息几何”}，那里需要核验信息矩阵的可逆性、步长和分布族的边界等条件。
